\originalchapter{22.02 2010 年期中卷}
题源: 2010年3月18日 Mid-Term Exam, 2012课程的历史练习材料. 官方题卷的题4通量为 $6\times10^{22}$ s$^{-1}$, 对应答案改印成 $6\times10^{21}$; 本章按题卷数值. 答案路线据官方整理, 有冲突处复算注明.
\begin{exercise}\label{ex:2202-e2010-1}
原题 1 (十组短问). (1) 可观测量及数学对象? (2) 测量可能值如何确定? (3) $\psi(x,y,z)$ 在 $(\widetilde x,\widetilde y,0)$ 附近 $\dd V=\dd x\dd y\dd z$ 的概率? (4) 量子体系演化随机吗? (5) $\psi(t_1)=a_1w_1+a_{91}w_{91}$, 测 $E_{91}$ 在 $t_1,2t_1$ 各概率? (6) $\int w_n^*\phi\dd x$ 是概率吗? (7) $\psi=B\cos(\sqrt2x)$ 的动量结果和概率? (8) $\psi=A\e^{-\ii79kx}$ 的动量、概率、物理意义? (9) $|l=5,m_z=-4\rangle$ 的 $L^2,L_z,L_x$ 可能值? (10) 两 $s=1/2,l=0$ 核子每个、总体各多少自旋态, 耦合表象标签?
\end{exercise}
\begin{solution}
(1) 可测量对应自伴算符. (2) 其谱中本征值或连续谱值. (3) $|\psi(\widetilde x,\widetilde y,0)|^2\dd V$. (4) 封闭体系 Schrödinger 演化确定且幺正; 测量结果有概率. (5) 若 H 不随时间变且无中途测量, 两时刻都 $|a_{91}|^2$. (6) 是振幅, 概率为模平方, 简并须投影子空间. (7) $p=\pm\sqrt2\hbar$, 各1/2, 在广义/盒归一意义下. (8) 唯一 $p=-79\hbar k$, 概率1, 是向负 x 的定动量自由态. (9) $L^2=30\hbar^2$, $L_z=-4\hbar$ 确定; $L_x$ 的谱为 $m_x\hbar$, $m_x=-5,\ldots,5$, 各概率可由旋转矩阵得, 不是全都相等. (10) 每个两态, 总四态; 标签 $s,m_s,s_1,s_2$, 分解为 s=0与1.
\end{solution}
\begin{exercise}\label{ex:2202-e2010-2}
原题 2. 成绩算符 $G\phi_k=g_k\phi_k$, 四正交归一态, $g=(A,B,C,D)$. 归一态 $\psi=c_1\phi_1+\phi_2/\sqrt5+2\phi_3/5+\phi_4/5$. (a) A 的概率? (b) 用 A=5,B=4,C=3,D=2 求均值. (c) Pass/Fail 算符在 k=1,2,3 时本征值 P, k=4 时 F, 与 G 对易吗, 说明.
\end{exercise}
\begin{solution}
(a) 归一化 $|c_1|^2=1-1/5-4/25-1/25=3/5$. (b) $\langle G\rangle=5(3/5)+4/5+12/25+2/25=109/25=4.36$. (c) 二者在同一基都对角, 逐态 $GP_F\phi_k=g_kp_k\phi_k=P_FG\phi_k$, 因而对易. P/F 的简并不妨碍该算符关系.
\end{solution}
\begin{exercise}\label{ex:2202-e2010-3}
原题 3. (a) 三段高度不同、含一个狭窄较高台阶的势中, 按下列源图性质将 A--D 波形同能量线1--4配对, 并解释波长、振幅、传播或衰减: A 在 I 振荡、跨高台阶衰减后中间低段又振荡、最右衰减; B 各段振荡; C 在 I、狭台阶、中间段振荡, 最右衰减; D 仅 I 振荡其后衰减. 能线1低于全部右侧台阶, 2高于狭台阶及中间段而低于最右台阶, 3高于全部势, 4高于中间段而低于狭台阶及最右段. (b) 两个三段波形 A、B 均传播, A 中间波长较长, B 中间较短; 匹配中央井 (势1) 与中央垒 (势2).
\end{exercise}
\begin{solution}
官方源图标签给 (a) A--4, B--3, C--2, D--1. 必须结合波形在最右段是否持续振荡核对: 源图 B 在全部可见段振荡并对应最高可传播能线3, C 在最右段衰减对应能线2, 其内部波数因两条高能线关系而不同. 上述源图无数值刻度, 本题以原图的相对波长和能线顺序为准, 不以图像的颜色区块当作精确势值. 允许区 $k=\sqrt{2m(E-V)}/\hbar$ 为实, 势较高处波长较长; 禁阻区有指数项. D 在所有右侧势都高于 E 时完全反射, A 可穿过较窄禁阻段进入后续允许段再遇最右禁阻区. (b) A--2, B--1, 中央垒降低动能、加长波长, 中央井反之.

本题图形的完整数值势值、幅度与尺度未给, 只进行定性配对. 原图与答案之间的台阶能线描述按页面位置核对, 下列可编辑表给各波形的关键判据, 不冒充精确本征函数:
\begin{center}\begin{tabular}{lll}\toprule
标签&判据&配对\\\midrule
A&窄禁阻段后又有允许区&4\\B&各区振荡,右段较长波&3\\C&前部各区振荡,右段衰减&2\\D&右侧只衰减&1\\\bottomrule\end{tabular}\end{center}
\end{solution}
\begin{exercise}\label{ex:2202-e2010-4}
原题 4. $\alpha$ 束通量 $\Gamma=6\times10^{22}$ s$^{-1}$, E=5 MeV, 方势垒 $V_B=85$ MeV, L=10 fm, $m_\alpha c^2=4000$ MeV, $\hbar c=200$ MeV fm. (a) 画波形. (b) 估计穿透概率. (c) 入射 $\psi=A\e^{\ii kx}$, 用 $\Gamma=|A|^2v$ 求 A. (d) 多久期望有一粒子透过?
\end{exercise}
\begin{solution}
(a) 左边为入射加反射, 势垒内两指数, 右边仅向右小幅透射, 左右波长同. (b) $\kappa=4$ fm$^{-1}$, 指数80. 正确厚方垒前因子为 $16E(V-E)/V^2=0.8858$, 因而 $T\simeq1.60\times10^{-35}$. 源答案用 $4\e^{-80}$ 作量级估算, 但 4 只在 E=V/2 适用, 本题两者相差约4.5倍. (c) $v=c\sqrt{2E/(mc^2)}=c/20=1.5\times10^7$ m/s, $|A|^2=4.0\times10^{15}\,\mathrm m^{-1}$, $|A|=6.32\times10^7\,\mathrm m^{-1/2}$, 任意整体相位. 此处波函数表示一维束粒子数密度而非单粒子归一概率; 通量单位据题意为每秒. (d) 期望等待时间 $1/(\Gamma T)\simeq1.04\times10^{12}$ s, 约 $3.3\times10^4$ 年. 若沿用源简略前因子4, 则为约 $7.3\times10^3$ 年; 必须同时声明前因子和通量版本.
\begin{center}\begin{tikzpicture}[x=1cm,y=.65cm]
\draw[->](-3,0)--(4,0)node[right]{$x$};\draw[->](-3,0)--(-3,3)node[above]{$V,\operatorname{Re}\psi$};
\draw[thick](-3,0)--(0,0)--(0,2.5)--(1.7,2.5)--(1.7,0)--(4,0);
\draw[dashed](-3,1.6)--(4,1.6)node[right]{$E<V$};
\draw[structurecolor,thick,domain=-3:0,samples=80]plot(\x,{.9*cos(180*\x)});
\draw[structurecolor,thick,domain=0:1.7,samples=50]plot(\x,{.9*exp(-2.2*\x)});
\draw[structurecolor,thick,domain=1.7:4,samples=80]plot(\x,{.022*cos(180*(\x-1.7))});
\node[below]at(-1.5,-1){入射与反射叠加};\node[below]at(.85,-1){指数区};\node[below]at(2.8,-1){小幅透射};
\draw[<->](0,2.9)--(1.7,2.9)node[midway,above]{$L$};
\end{tikzpicture}\end{center}
图仅显示振荡、指数和透射幅度的定性关系, 各段示意曲线不用于读取匹配常数; 精确复波函数应同时匹配函数与导数.

\end{solution}
